\documentclass{article}
\usepackage[dvipsnames,svgnames,x11names]{xcolor}
\usepackage{tikz}
\usepackage{amsmath}
\usepackage{amssymb}
\usepackage{geometry}
\usepackage{subcaption}
\pagecolor{white}

\usetikzlibrary{backgrounds}
\usetikzlibrary{arrows.meta}
\usetikzlibrary{fit}

\begin{document}

\section*{TIDBIT DUMP {\em (inexhaustive)}}

\begin{itemize}
\item[:)]{An, ty A. Basically, GR extend is assert all subspace curvature is extrinsic, while base manifold is intrinsically flat, though perhaps subject to the 'pullback metric' thingy.}
\end{itemize}

\section*{A publication by the Time Machine For Peace Social Invention Program.}

\noindent Define, basis:\\
$\hat e = (\hat e_1, \hat e_2, \hat e_3, \hat e_4, \hat e_5, \hat e_6, \hat e_7)$\\
$\hat e = (\hat\tau, \hat x, \hat t_x, \hat y, \hat t_y, \hat z, \hat t_z)$\\

\\
Vector tuples:\\
$\vec U = (u_\tau, u_x, u_{t_x}, u_y, u_{t_y}, u_z, u_{t_z})$\\
$\vec V = (v_\tau, v_x, v_{t_x}, v_y, v_{t_y}, v_z, v_{t_z})$\\
$\vec W = (w_\tau, w_x, w_{t_x}, w_y, w_{t_y}, w_z, w_{t_z})$ 


\section*{Dot Product in $\mathbb{R}^7$:}
\begin{align*}
W = U \cdot V =& u_\tau v_\tau \hat \tau +u_x v_x \hat x + u_{t_x} v_{t_x} \hat t_x + u_y v_y \hat y + u_{t_y} v_{t_y} \hat t_y + u_z v_z \hat z + u_{t_z} v_{t_z} \hat t_z\\
=& (u_\tau v_\tau, u_x v_x, u_{t_x} v_{t_x}, u_y v_y, u_{t_y} v_{t_y}, u_z v_z, u_{t_z} v_{t_z})\\
\end{align*}

\section*{Cross Product in $\mathbb{R}^7$:}
\begin{align*}
U(\times_{7}) =
\begin{pmatrix}
0 & -u_{t_x} & u_x & -u_{t_y} & u_y & u_{t_z} & -u_z \\[1mm]
u_{t_x} & 0 & -u_\tau & -u_z & -u_{t_z} & u_y & u_{t_y} \\[1mm]
-u_x & u_\tau & 0 & -u_{t_z} & u_z & -u_{t_y} & u_y \\[1mm]
u_{t_y} & u_z & u_{t_z} & 0 & -u_\tau & -u_x & -u_{t_x} \\[1mm]
-u_y & u_{t_z} & -u_z & u_\tau & 0 & u_{t_x} & -u_x \\[1mm]
-u_{t_z} & -u_y & u_{t_y} & u_x & -u_{t_x} & 0 & u_\tau \\[1mm]
u_z & -u_{t_y} & -u_y & u_{t_x} & u_x & -u_\tau & 0 \\[1mm]
\end{pmatrix}
\end{align*}

\setlength{\jot}{7pt}
\begin{align*}
W = U(\times_{7})V =&
\begin{pmatrix}
0 & -u_{t_x} & u_x & -u_{t_y} & u_y & u_{t_z} & -u_z \\[1mm]
u_{t_x} & 0 & -u_\tau & -u_z & -u_{t_z} & u_y & u_{t_y} \\[1mm]
-u_x & u_\tau & 0 & -u_{t_z} & u_z & -u_{t_y} & u_y \\[1mm]
u_{t_y} & u_z & u_{t_z} & 0 & -u_\tau & -u_x & -u_{t_x} \\[1mm]
-u_y & u_{t_z} & -u_z & u_\tau & 0 & u_{t_x} & -u_x \\[1mm]
-u_{t_z} & -u_y & u_{t_y} & u_x & -u_{t_x} & 0 & u_\tau \\[1mm]
u_z & -u_{t_y} & -u_y & u_{t_x} & u_x & -u_\tau & 0 \\[1mm]
\end{pmatrix}
\begin{pmatrix}
v_\tau  \\[1mm]
v_x \\[1mm]
v_{t_x} \\[1mm]
v_y \\[1mm]
v_{t_y} \\[1mm]
v_z \\[1mm]
v_{t_z} \\[1mm]
\end{pmatrix}\\[3mm]
\
=& \quad (u_x v_{t_x} - u_{t_x} v_x  + u_y v_{t_y} - u_{t_y} v_y + u_{t_z} v_z  - u_z v_{t_z}) \hat \tau \\
+& \quad (u_{t_x} v_\tau - u_\tau v_{t_x} + u_y v_z - u_z v_y  + u_{t_y} v_{t_z} - u_{t_z} v_{t_y}) \hat x \\
+& \quad (u_\tau v_x - u_x v_\tau + u_y v_{t_z} - u_{t_z} v_y + u_z v_{t_y} - u_{t_y} v_z) \hat t_x \\
+& \quad (u_{t_y} v_\tau - u_\tau v_{t_y} + u_z v_x - u_x v_z + u_{t_z} v_{t_x} - u_{t_x} v_{t_z}) \hat y \\  \tag{FORM 1}
+& \quad (u_\tau v_y - u_y v_\tau + u_{t_x} v_z - u_z v_{t_x} + u_{t_z} v_x - u_x v_{t_z}) \hat t_y \\
+& \quad (u_\tau v_{t_z} - u_{t_z} v_\tau + u_x v_y - u_y v_x + u_{t_y} v_{t_x} - u_{t_x} v_{t_y}) \hat z \\
+& \quad (u_z v_\tau - u_\tau v_z + u_x v_{t_y} - u_{t_y} v_x + u_{t_x} v_y - u_y v_{t_x}) \hat t_z \\[3mm]
=& \quad (u_{[x}v_{t_x]}  + u_{[y} v_{t_y]} + u_{[t_z} v_{z]}) \hat \tau \\
+& \quad (u_{[t_x} v_{\tau]} + u_{[y} v_{z]} + u_{[t_y} v_{t_z]}) \hat x \\
+& \quad (u_{[\tau} v_{x]} + u_{[y} v_{t_z]} + u_{[z} v_{t_y]}) \hat t_x \\
+& \quad (u_{[t_y} v_{\tau]} + u_{[z} v_{x]} + u_{[t_z} v_{t_x]}) \hat y \\ \tag{FORM 2}
+& \quad (u_{[\tau} v_{y]} + u_{[t_x} v_{z]} + u_{[t_z} v_{x]}) \hat t_y \\
+& \quad (u_{[\tau} v_{t_z]} + u_{[x} v_{y]} + u_{[t_y} v_{t_x]}) \hat z \\
+& \quad (u_{[z} v_{\tau]} + u_{[x} v_{t_y]} + u_{[t_x} v_{y]}) \hat t_z \\
=& \quad (u_x v_{t_x} - u_{t_x} v_x  + u_y v_{t_y} - u_{t_y} v_y + u_{t_z} v_z  - u_z v_{t_z}) \hat \tau \\
+& \quad v_\tau[(u_{t_x}\hat x - u_x \hat t_x) + (u_{t_y}\hat y - u_y \hat t_y) - (u_{t_z}\hat z - u_z \hat t_z)]\\
-& \quad u_\tau[(v_{t_x}\hat x - v_x \hat t_x) + (v_{t_y}\hat y - v_y \hat t_y) - (v_{t_z}\hat z - v_z \hat t_z)]\\
+& \quad v_{t_z}[(u_{t_y}\hat x + u_y \hat t_x) - (u_{t_x}\hat y + u_x \hat t_y)]\\
-& \quad u_{t_z}[(v_{t_y}\hat x + v_y \hat t_x) - (v_{t_x}\hat y + v_x \hat t_y)]\\
+& \quad v_z[(u_y\hat x - u_{t_y} \hat t_x) - (u_x\hat y - u_{t_x} \hat t_y)]\\
-& \quad u_z[(v_y\hat x - v_{t_y} \hat t_x) - (v_x\hat y - v_{t_x} \hat t_y)]\\  \tag{FORM 3}
+& \quad v_{t_x}(u_{t_y}\hat z - u_y \hat t_z) - v_x(u_y\hat z + u_{t_y} \hat t_z)\\
-& \quad u_{t_x}(v_{t_y}\hat z - v_y \hat t_z) + u_x(v_y\hat z + v_{t_y} \hat t_z)\\
\end{align*}

\section*{Gradient in $\mathbb{R}^7$:}

\noindent Define scalar field, $\phi: \mathbb{R}^7 \to \mathbb{R}^1$:\\

\noindent Define del operator $\nabla_7$, in context of vector field $F$:\\
\\
$\nabla_7 = \frac{\partial}{\partial \tau}\hat{\tau} + \frac{\partial}{\partial x}\hat{x} + \frac{\partial}{\partial t_x}\hat{t}_x + \frac{\partial}{\partial y}\hat{y} + \frac{\partial}{\partial t_y}\hat{t}_y + \frac{\partial}{\partial z}\hat{z} + \frac{\partial}{\partial t_z}\hat{t}_z$\\
\\
So,\\
\\
$\nabla_7 \phi = \frac{\partial \phi}{\partial \tau}\hat{\tau} + \frac{\partial \phi}{\partial x}\hat{x} + \frac{\partial \phi}{\partial t_x}\hat{t}_x + \frac{\partial \phi}{\partial y}\hat{y} + \frac{\partial \phi}{\partial t_y}\hat{t}_y + \frac{\partial \phi}{\partial z}\hat{z} + \frac{\partial \phi}{\partial t_z}\hat{t}_z$\\

\section*{Divergence in $\mathbb{R}^7$:}

$\nabla_7 \cdot F = \frac{\partial}{\partial \tau} F_\tau + \frac{\partial}{\partial x} F_x + \frac{\partial}{\partial t_x} F_{t_x} + \frac{\partial}{\partial y} F_y + \frac{\partial}{\partial t_y} F{t_y} + \frac{\partial}{\partial z} F_z + \frac{\partial}{\partial t_z} F_{t_z}$\\

\section*{Curl in $\mathbb{R}^7$:}


\noindent Define vector field, $F: \mathbb{R}^7 \to \mathbb{R}^7$:\\
\\
$F = F_\tau \hat{\tau} + F_x \hat{x} + F_{t_x} \hat{t}_x + F_y \hat{y} + F_{t_y} \hat{t}_y + F_z \hat{z} + F_{t_z} \hat{t}_z$\\




The curl of $F$ in $\mathbb{R}^7$, denoted $\nabla \times_{7} F$, is defined as the 7-dimensional cross product of the del operator $\nabla$ with the vector field $F$:

\setlength{\jot}{7pt}
\begin{align*}
\nabla \times_{7} F = \nabla(\times_{7})F =& \left(\frac{\partial F_{t_x}}{\partial x} - \frac{\partial F_x}{\partial t_x} + \frac{\partial F_{t_y}}{\partial y} - \frac{\partial F_y}{\partial t_y} + \frac{\partial F_z}{\partial t_z} - \frac{\partial F_{t_z}}{\partial z}\right) \hat{\tau} \\
+& \left(\frac{\partial F_\tau}{\partial t_x} - \frac{\partial F_{t_x}}{\partial \tau} + \frac{\partial F_z}{\partial y} - \frac{\partial F_y}{\partial z} + \frac{\partial F_{t_z}}{\partial t_y} - \frac{\partial F_{t_y}}{\partial t_z}\right) \hat{x} \\
+& \left(\frac{\partial F_x}{\partial \tau} - \frac{\partial F_\tau}{\partial x} + \frac{\partial F_{t_z}}{\partial y} - \frac{\partial F_y}{\partial t_z} + \frac{\partial F_{t_y}}{\partial z} - \frac{\partial F_z}{\partial t_y}\right) \hat{t}_x \\
+& \left(\frac{\partial F_\tau}{\partial t_y} - \frac{\partial F_{t_y}}{\partial \tau} + \frac{\partial F_x}{\partial z} - \frac{\partial F_z}{\partial x} + \frac{\partial F_{t_x}}{\partial t_z} - \frac{\partial F_{t_z}}{\partial t_x}\right) \hat{y} \\
+& \left(\frac{\partial F_y}{\partial \tau} - \frac{\partial F_\tau}{\partial y} + \frac{\partial F_z}{\partial t_x} - \frac{\partial F_{t_x}}{\partial z} + \frac{\partial F_x}{\partial t_z} - \frac{\partial F_{t_z}}{\partial x}\right) \hat{t}_y \\
+& \left(\frac{\partial F_{t_z}}{\partial \tau} - \frac{\partial F_\tau}{\partial t_z} + \frac{\partial F_y}{\partial x} - \frac{\partial F_x}{\partial y} + \frac{\partial F_{t_x}}{\partial t_y} - \frac{\partial F_{t_y}}{\partial t_x}\right) \hat{z} \\
+& \left(\frac{\partial F_z}{\partial \tau} - \frac{\partial F_\tau}{\partial z} + \frac{\partial F_{t_y}}{\partial x} - \frac{\partial F_x}{\partial t_y} + \frac{\partial F_y}{\partial t_x} - \frac{\partial F_{t_x}}{\partial y}\right) \hat{t}_z\\[3mm]
=& \left(\partial_{[x}F_{t_x]} + \partial_{[y}F_{t_y]} + \partial_{[t_z}F_{z]}\right) \hat{\tau} \\
+& \left(\partial_{[t_x}F_{\tau]} + \partial_{[y}F_{z]} + \partial_{[t_y}F_{t_z]}\right) \hat{x} \\
+& \left(\partial_{[\tau}F_{x]} + \partial_{[y}F_{t_z]} + \partial_{[z}F_{t_y]}\right) \hat{t}_x \\
+& \left(\partial_{[t_y}F_{\tau]} + \partial_{[z}F_{x]} + \partial_{[t_z}F_{t_x]}\right) \hat{y} \\
+& \left(\partial_{[\tau}F_{y]} + \partial_{[t_x}F_{z]} + \partial_{[t_z}F_{x]}\right) \hat{t}_y \\
+& \left(\partial_{[\tau}F_{t_z]} + \partial_{[x}F_{y]} + \partial_{[t_y}F_{t_x]}\right) \hat{z} \\
+& \left(\partial_{[z}F_{\tau]} + \partial_{[x}F_{t_y]} + \partial_{[t_x}F_{y]}\right) \hat{t}_z
\end{align*}

\section*{Curl Properties in $\mathbb{R}^7$}
\subsection*{$\nabla (\times_7) (\nabla_7 \phi) = 0$ for any scalar field $\phi$}
Let $\phi$ be a scalar field in $\mathbb{R}^7$. The gradient of $\phi$ is:
\begin{align*}
\nabla \phi = \frac{\partial \phi}{\partial \tau}\hat{\tau} + \frac{\partial \phi}{\partial x}\hat{x} + \frac{\partial \phi}{\partial t_x}\hat{t}_x + \frac{\partial \phi}{\partial y}\hat{y} + \frac{\partial \phi}{\partial t_y}\hat{t}_y + \frac{\partial \phi}{\partial z}\hat{z} + \frac{\partial \phi}{\partial t_z}\hat{t}_z
\end{align*}
To compute $\nabla \times_7 (\nabla \phi)$, we use the component form of the curl. For the $\hat{\tau}$ component:
\begin{align*}
[\nabla \times_7 (\nabla \phi)]_\tau &= \frac{\partial}{\partial x}\left(\frac{\partial \phi}{\partial t_x}\right) - \frac{\partial}{\partial t_x}\left(\frac{\partial \phi}{\partial x}\right) \
&+ \frac{\partial}{\partial y}\left(\frac{\partial \phi}{\partial t_y}\right) - \frac{\partial}{\partial t_y}\left(\frac{\partial \phi}{\partial y}\right) \
&+ \frac{\partial}{\partial t_z}\left(\frac{\partial \phi}{\partial z}\right) - \frac{\partial}{\partial z}\left(\frac{\partial \phi}{\partial t_z}\right)
\end{align*}
By Schwarz's theorem on the equality of mixed partial derivatives (assuming $\phi$ is $C^2$), each pair cancels out:
\begin{align*}
\frac{\partial}{\partial x}\left(\frac{\partial \phi}{\partial t_x}\right) = \frac{\partial^2 \phi}{\partial x \partial t_x} = \frac{\partial^2 \phi}{\partial t_x \partial x} = \frac{\partial}{\partial t_x}\left(\frac{\partial \phi}{\partial x}\right)
\end{align*}
Similarly, all other pairs in each component will cancel, giving us:
\begin{align*}
\nabla \times_7 (\nabla \phi) = \vec{0}
\end{align*}

\subsection{R3 Subspace Interpretation}
Each component of the curl can be interpreted as the sum of three terms, where each term corresponds to the curl operation in a specific R3 subspace of R7:

\begin{align*}
(\nabla \times_{7} F)_\tau &= \text{curl in } (x,t_x) \text{ plane} + \text{curl in } (y,t_y) \text{ plane} + \text{curl in } (z,t_z) \text{ plane}\\
(\nabla \times_{7} F)_x &= \text{curl in } (\tau,t_x) \text{ plane} + \text{curl in } (y,z) \text{ plane} + \text{curl in } (t_y,t_z) \text{ plane}\\
\end{align*}
and so on for the other components.

\subsection{Properties of the Curl in R7}
The curl in R7 maintains several important properties analogous to the R3 curl:

1. $\nabla \times_{7} (\nabla \phi) = 0$ for any scalar field $\phi$

2. $\nabla \cdot (\nabla \times_{7} F) = 0$ for any vector field $F$

3. The curl measures the infinitesimal rotation of the vector field around each axis.

However, due to the non-associativity of the 7D cross product, some identities from R3 vector calculus need to be modified appropriately when working in R7.


\section{Demonstration of Key Curl Properties in $\mathbb{R}^7$}
\subsection{Property 1: $\nabla \times_7 (\nabla \phi) = 0$ for any scalar field $\phi$}
Again, $\phi$ a scalar field in $\mathbb{R}^7$. Gradient of $\phi$:
\begin{align*}
\nabla \phi = \frac{\partial \phi}{\partial \tau}\hat{\tau} + \frac{\partial \phi}{\partial x}\hat{x} + \frac{\partial \phi}{\partial t_x}\hat{t}_x + \frac{\partial \phi}{\partial y}\hat{y} + \frac{\partial \phi}{\partial t_y}\hat{t}_y + \frac{\partial \phi}{\partial z}\hat{z} + \frac{\partial \phi}{\partial t_z}\hat{t}_z
\end{align*}
To compute $\nabla \times_7 (\nabla \phi)$, we use the component form of the curl. For the $\hat{\tau}$ component:
\begin{align*}
[\nabla \times_7 (\nabla \phi)]_\tau &= \frac{\partial}{\partial x}\left(\frac{\partial \phi}{\partial t_x}\right) - \frac{\partial}{\partial t_x}\left(\frac{\partial \phi}{\partial x}\right) \
&+ \frac{\partial}{\partial y}\left(\frac{\partial \phi}{\partial t_y}\right) - \frac{\partial}{\partial t_y}\left(\frac{\partial \phi}{\partial y}\right) \
&+ \frac{\partial}{\partial t_z}\left(\frac{\partial \phi}{\partial z}\right) - \frac{\partial}{\partial z}\left(\frac{\partial \phi}{\partial t_z}\right)
\end{align*}
By Schwarz's theorem on the equality of mixed partial derivatives (assuming $\phi$ is $C^2$), each pair cancels out:
\begin{align*}
\frac{\partial}{\partial x}\left(\frac{\partial \phi}{\partial t_x}\right) = \frac{\partial^2 \phi}{\partial x \partial t_x} = \frac{\partial^2 \phi}{\partial t_x \partial x} = \frac{\partial}{\partial t_x}\left(\frac{\partial \phi}{\partial x}\right)
\end{align*}
Similarly, all other pairs in each component will cancel, giving us:
\begin{align*}
\nabla \times_7 (\nabla \phi) = \vec{0}
\end{align*}
\subsection{Property 2: $\nabla \cdot (\nabla \times_7 F) = 0$ for any vector field $F$}
For a vector field $F$ in $\mathbb{R}^7$, let's compute $\nabla \cdot (\nabla \times_7 F)$:
\begin{align*}
\nabla \cdot (\nabla \times_7 F) = \frac{\partial}{\partial \tau}[\nabla \times_7 F]_\tau + \frac{\partial}{\partial x}[\nabla \times_7 F]x + \frac{\partial}{\partial t_x}[\nabla \times_7 F]{t_x} + \ldots
\end{align*}
Substituting the components of $\nabla \times_7 F$:
\begin{align*}
\nabla \cdot (\nabla \times_7 F) &= \frac{\partial}{\partial \tau}\left(\frac{\partial F_{t_x}}{\partial x} - \frac{\partial F_x}{\partial t_x} + \frac{\partial F_{t_y}}{\partial y} - \frac{\partial F_y}{\partial t_y} + \frac{\partial F_z}{\partial t_z} - \frac{\partial F_{t_z}}{\partial z}\right) \
&+ \frac{\partial}{\partial x}\left(\frac{\partial F_\tau}{\partial t_x} - \frac{\partial F_{t_x}}{\partial \tau} + \frac{\partial F_z}{\partial y} - \frac{\partial F_y}{\partial z} + \frac{\partial F_{t_z}}{\partial t_y} - \frac{\partial F_{t_y}}{\partial t_z}\right) \
&+ \frac{\partial}{\partial t_x}\left(\frac{\partial F_x}{\partial \tau} - \frac{\partial F_\tau}{\partial x} + \frac{\partial F_{t_z}}{\partial y} - \frac{\partial F_y}{\partial t_z} + \frac{\partial F_{t_y}}{\partial z} - \frac{\partial F_z}{\partial t_y}\right) \
&+ \ldots
\end{align*}
When we expand this expression fully, we get terms like:
\begin{align*}
\frac{\partial^2 F_{t_x}}{\partial \tau \partial x} - \frac{\partial^2 F_x}{\partial \tau \partial t_x} + \ldots
\end{align*}
For each term $\frac{\partial^2 F_a}{\partial b \partial c}$ that appears, there will be exactly one corresponding term $-\frac{\partial^2 F_a}{\partial b \partial c}$ elsewhere in the expansion.
For example, the term $\frac{\partial^2 F_{t_x}}{\partial \tau \partial x}$ from the $\tau$-component is canceled by $-\frac{\partial^2 F_{t_x}}{\partial \tau \partial x}$ from the $x$-component.
This systematic cancellation occurs because of the antisymmetric structure of the curl operator and the equality of mixed partial derivatives. When fully expanded, all terms cancel in pairs, giving:
\begin{align*}
\nabla \cdot (\nabla \times_7 F) = 0
\end{align*}


\section*{Rotation matrix}

\begin{align}
\times_{7R} =
\begin{pmatrix}
R_{\tau\tau} & R_{0S} \\[1mm]
R_{S0} & R_{SS}
\end{pmatrix}
\end{align}

\begin{align}
R_{SS} =
\begin{pmatrix}
R^{(x)} &R^{(x,y)} & R^{(x,z)} \\
R^{(y,x)} &R^{(y)} & R^{(y,z)} \\
R^{(z,x)} &R^{(z,y)} & R^{(z)}
\end{pmatrix}
\end{align}



\newcommand{\threedaxes}[7][1]{%
    % #1 = scale (optional, default 0.8)
    % #2 = i-axis label
    % #3 = i-axis length
    % #4 = j-axis label
    % #5 = j-axis length
    % #6 = l=axis label
    % #7 = k-axis direction (1 for out, -1 for in)
    \begin{tikzpicture}[scale=#1]
        % Define axis style
        \tikzset{axis/.style={-{Stealth[length=#1mm, width=#1mm]}, thick}}
        % i and j axes
        \draw[axis] (0,0) -- (#3,0) node at (#3-#3/10,-0.2) {\small $#2$};
        \draw[axis] (0,0) -- (0,#5) node at (-0.3,#5-#5/20) {\small $#4$};
        
        % Draw a white-filled circle first to cover the axis lines
        \fill[white] (0,0) circle (0.12);
        % k axis 
        \draw[thick] (0,0) circle (0.1);
        \ifnum#7<0
            % Into of page (cross)
            \draw[thick] (0,0) -- (0.07,0.07);
            \draw[thick] (0,0) -- (-0.07,-0.07);
            \draw[thick] (0,0) -- (0.07,-0.07);
            \draw[thick] (0,0) -- (-0.07,0.07);
            \node at (0,-0.25) {\small $#6$};
        \else
        \ifnum#7=0
            % Out of page (dot)
            \draw[thick] (0,0) -- (0.07,0.07);
            \draw[thick] (0,0) -- (-0.07,-0.07);
            \draw[thick] (0,0) -- (0.07,-0.07);
            \draw[thick] (0,0) -- (-0.07,0.07);
            \fill[black] (0,0) circle (0.05);
            \node at (0,-0.25) {\small $#6$};
        \else
            % Out of page (dot)
            \fill[black] (0,0) circle (0.05);
            \node at (0,-0.25) {\small $#6$};
        \fi
        \fi
    \end{tikzpicture}%
}

\clearpage
\begin{figure}
\begin{minipage}{0.3\textwidth}
	\centering
	\threedaxes[2]{\{i, t_i\}}{2}{\{t_i, i\}}{2}{\{\tau\}(\pm\pm\mp)^*}{0}
\end{minipage}
\hfill
\begin{minipage}{0.3\textwidth}
	\centering
	\threedaxes[2]{\{\tau, i\}}{2}{\{i, \tau\}}{2}{\{t_i\}(\pm\pm\mp)^*}{0}
\end{minipage}
\hfill
\begin{minipage}{0.3\textwidth}
	\centering
	\threedaxes[2]{\{\tau, t_i\}}{2}{\{t_i, \tau\}}{2}{\{i\}(\mp\mp\pm)^*}{0}
\end{minipage}

\vspace{1cm}

\begin{minipage}{0.3\textwidth}
	\centering
	\threedaxes[2]{\{t_i, i\}}{2}{\{t_j, j\}}{2}{\{k\}(\mp++)^*}{0}
\end{minipage}

\vspace{1cm}

\begin{minipage}{0.3\textwidth}
	\centering
	\threedaxes[2]{\{t_i, i\}}{2}{\{j, t_j\}}{2}{\{t_k\}(+\mp\pm)^*}{0}
\end{minipage}
\end{figure}




TIME MACHINE FOR PEACE SOCIAL INVENTION PROGRAM
\end{document}